% 【本文件写什么】独立子 crate，无 api/impl 目录树：base-config
% - 职责：base 配置子 crate，无 api/impl 分叉。
% - 源码：os/components/wateros-base/base-config/src/
% 【事实来源】
% - os/components/wateros-base/base-config/、父组件 src/lib.rs

\subsection{base-config}

\paragraph{wateros-base 类型层}

聚合\code{src/lib.rs}导出四模块：
\begin{itemize}
 \item \code{addr}：\code{BasePhysAddr}、\code{BaseVirtAddr}、\code{BasePPN}、\code{BaseVPN}，均为 \code{repr(C)}，字段 \code{val: usize}。\code{BasePhysAddr}/\code{BaseVirtAddr}实现\code{Into<*mut T>}薄转换，不校验对齐或映射有效性；
 \item \code{boot}：\code{DTBPA = BasePhysAddr}，设备树物理基址类型别名；
 \item \code{cpu}：\code{CPUHartID = usize}；
 \item \code{sync}：\code{UniprocessorSafeCell<T>}，单核\code{RefCell}包装，\code{unsafe fn new}与\code{exclusive\_access() -> RefMut}。多hart或抢占并发须用平台锁，勿误用此容器。
\end{itemize}

\paragraph{wateros-base-config 常量层}

独立crate\code{wateros-base-config}，\code{\#![no\_std]}，按子系统分模块：

\code{syscall}：\code{MAX\_SYSCALL\_ARGS = 6}，与陷阱帧及\code{SyscallArgs}数组长度一致。

\code{mm}：内核堆\code{KERNEL\_HEAP\_SIZE = 1 << 27}，128MiB；QEMU \code{virt}物理RAM区间\code{0x8000\_0000}--\code{0xC000\_0000}；MMIO半开区间\code{0x1000\_0000}--\code{0x1200\_0000}。真机或自定义\code{-m}应以DTB为准，此处为bring-up回退。

\code{ipc}：\code{DEFAULT\_PIPE\_CAPACITY = 4096}。

\code{fs}：页缓存\code{FILE\_PAGE\_SIZE = 4096}、大文件阈值64KiB、全局页缓存4096槽、读预取步长8页、\code{FileIoMode::Direct}，\code{Async} 在配置层存在但 v1 未实现，块设备LBA缓存1024块。

\code{task}：定时器周期10ms、普通任务最大50 tick、RR任务最大10 tick、ready队列stale compact阈值8。

\code{klog}：\code{KLOG\_DESC\_SLOTS = 256}、\code{KLOG\_TEXT\_RING\_BYTES = 32 * 1024}、\code{KLOG\_MAX\_RECORD\_BYTES = 1024}。

\paragraph{设计边界}

\code{wateros-base}只放类型，不放配置数值；配置数值只在\code{wateros-base-config}维护。地址newtype不附带地址空间标识校验；\code{UniprocessorSafeCell}仅适用于单核假设。
